% 单位圆面积相位与相位纤维 T 的接口：终点断言与实轴具体化
% 基于正文三条可审计接口（可逆演化—部分观测—扩展保存；相位纤维 T；单位圆扭曲下的相位–幅度分离）
% 给出与全文投影词/完成化语言兼容的终点断言簇，不复述既有编号结论。

\subsection{单位圆面积相位与相位纤维 \(\mathcal{T}\) 的接口：终点断言与实轴具体化}\label{subsec:unit-circle-phase-T-interface}
在正文已确立的三条可审计接口——(i) 可逆演化、部分观测与扩展保存（信息守恒严格化为扩展空间上的单射/可逆性）；(ii) 相位纤维 \(\mathcal{T}\) 作为紧阿贝尔群及其角色基底，将可加观测量提升为可乘权重并实现卷积对角化（定义 \ref{def:pom-phase-fiber}）；(iii) 单位圆扭曲下的相位–幅度分离（压力的相位项线性、幅度项仅依赖偶对称变量，见命题 \ref{prop:completion-endpoint-jet-diffusion} 与推论 \ref{cor:sync-kernel-weighted-phase-amplitude}）——之下，单位圆面积相位算术可被提升为与全文投影词及完成化语言兼容、且可产生终点断言的结构层。以下给出面积相位坐标与相位测度的规范表述，并给出一组可直接纳入正文“可审计输出”的终点断言与实轴具体化段落；各断言均以“最终陈述 + 可检验含义/可插入正文的接口形式”给出，且不重复既有编号结论。

\paragraph{面积相位坐标与相位测度（规范表述）}
令面积相位坐标由无参数积分确定：
\(\theta(x):=2\int_{0}^{x}(1+t^{2})^{-1}\dd t=2\arctan x\in(-\pi,\pi)\)，
并令 \(u(x):=e^{\mathrm{i}\theta(x)}\in\TT\)。
与 \S\ref{subsec:unit-circle-phase-lift} 的记号一致：\(\theta(x)\) 即式 \eqref{eq:unit-circle-phase-vartheta} 的 \(\vartheta(x)\)，\(u(x)\) 即式 \eqref{eq:unit-circle-cayley-embed} 的 \(\iota(x)\)。
以单位圆均匀测度 \(\dd\theta/(2\pi)\) 为基准，将其拉回实轴得到相位测度
\(\dd\mu(x)=(2\pi)^{-1}\dd\theta(x)=\bigl(\pi(1+x^{2})\bigr)^{-1}\dd x\)（与 \eqref{eq:unit-circle-phase-cauchy-measure} 一致）。
该测度实现“离零点越近权重越大”的严格化：相位坐标上为等距（均匀），实轴上密度为 \((1+x^{2})^{-1}\) 型衰减。

\subsubsection{终点断言与可检验含义}

\begin{proposition}[射影直线的单位圆规约与投影门同构]\label{prop:phase-projection-gate-isomorphism}
将实数视为射影直线 \(\mathbb{P}^{1}(\RR)\) 的非无穷点并施加单位规范（gauge fixing），则 \(u(x)=e^{\mathrm{i}\theta(x)}\) 给出一个无参数的规范截面，其结构等价于：先将对象提升到齐次坐标（过程层），再经单位圆规范化将结果投影回唯一代表（读出层）。
该“单位圆规范化”在逻辑类型上与正文“沿群胚轨道投影到横截面”的折叠算子同型：两者均属于可逆演化与末端一次投影门范式。
\end{proposition}

\begin{corollary}[相位投影门的常数预算]\label{cor:phase-projection-gate-budget}
对任意由相位算术基本运算复合得到的表达式，存在等价的实现使得单位圆规范化（相位投影门）在整条计算链中至多调用一次（仅在末端读出阶段）。
其语义地位与正文 \(\Fold_\infty\) 末端一次调用的预算结论同型（见 \S\ref{subsec:pom-state-evolution} 与 \ref{axiom:unit-circle-phase-extension-preserving}）。
\end{corollary}
\begin{remark}[可检验含义]
若在实现中观察到因“中途反复规范化”导致的数值或信息漂移，则该漂移应可被重写系统消去；若漂移不可消去，则表明实现中隐含了额外的序门或额外记录轴。
\end{remark}

\begin{proposition}[相位压缩的可逆化与残差轴外置]\label{prop:phase-residual-axis}
任何将非紧对象（如 \(\RR\) 或更高维实/复空间）经相位坐标压缩到紧相位纤维 \(\TT\) 的方案，若允许“提升—投影”反复出现，则不可避免地产生被投影抹去的尺度残差；若要求整体可逆，则残差必须外置为额外正交轴。
该轴的增长并非“无穷远处含无穷信息”，而是在扩展空间上新增正交自由度以保持单射，与正文可逆化接口（公理 \ref{axiom:unit-circle-phase-extension-preserving}）一致。
\end{proposition}
\begin{remark}[可检验含义]
在投影词语言下，所有“不可解释的误差/不可回收的溢出”应可被定位为某一步将残差未外置而直接丢弃，从而破坏 \(\widetilde\pi\) 的单射性。
\end{remark}

\begin{proposition}[紧相位塔与极限表示]\label{prop:compact-phase-tower}
将“提升”解释为把对象嵌入某个紧相位纤维（球面/圆环/高维环面等）并引入相应角色基底，则迭代提升产生一条由紧群组成的逆系统（或归纳系统）；其极限的 \(L^{2}\) 表示自然分裂为可分 Hilbert 张量因子，每一层新增的相位自由度在谱上对应新的正交模态。
相位纤维 \(\mathcal{T}\) 因而可解释为：紧群角色分解所必然带来的正交分解（定义 \ref{def:pom-phase-fiber}）。
\end{proposition}
\begin{remark}[可检验含义]
若相位提升层数增加而观测谱（非平凡角色的衰减/混合时间）不产生新的可分辨模态，则所谓“提升”并未真正增加正交自由度，而仅为重参数化（可被证伪）。
\end{remark}

\begin{proposition}[内禀扭曲原理]\label{prop:phase-intrinsic-twist}
若采用面积相位坐标，则“对核施加单位圆扭曲并研究其 Perron 根的相位–幅度分离”可被重新解释为：研究算术对象在相位纤维 \(\mathcal{T}\) 中的内禀位置变量对谱的响应。
扭曲角不再作为外参扫描，而是数轴在紧相位上的坐标；正文的相位纤维 \(\mathcal{T}\) 与完成化核的单位圆扭曲理论在语义上闭合为同一对象——相位坐标化后的算术状态空间（参见注 \ref{rem:unit-circle-phase-completion-half-angle}）。
\end{proposition}

\begin{proposition}[无穷问题的模态预算重表述]\label{prop:phase-infinity-modal-budget}
将非紧尺度压缩到紧相位后，“无穷远处”的技术困难等价于：为维持解析/谱统计的正确密度，非平凡角色的有效模态数必须随尺度增长；若将相位模态数固定在有限维（有限状态核/有限矩阵口径），则临界线几何不可避免地退化为有限条竖线与有限相位梳齿，从而无法承载 \(T\log T\) 级的统计结构。
该重表述将“无穷问题”从空间端点转移为扩展塔的不可省略性，与正文信息守恒框架一致（注 \ref{rem:app-rvm-endpoint-constant}）。
\end{proposition}

\begin{proposition}[相位算术的酉化归约]\label{prop:phase-unitary-reduction}
若将“数”的主语义全部放置在紧相位纤维上（即把主要不变量写成相位读出），则任一 RH 型“临界线”断言均可在相位侧被等价改写为：某个由投影词编译得到的伴随矩阵/转移算子在非平凡角色子空间上的酉性或拟酉性约束；所谓临界线相位即为该酉谱的相位。
面积相位坐标提供将外参扭曲内禀化的机制（命题 \ref{prop:phase-intrinsic-twist}），从而使酉化可被解释为算术在相位纤维上的闭合性条件（与命题 \ref{prop:xi-hilbert-polya-cmv} 等单位圆根–酉矩阵接口对齐）。
\end{proposition}

\paragraph{实轴具体化段落（可并入正文的接口表述）}
采用无参数积分 \(\theta(x)=2\int_0^x (1+t^2)^{-1}\dd t\) 定义实轴的面积相位坐标 \(u(x)=e^{\mathrm{i}\theta(x)}\in\TT\)。
其拉回测度为 \(\dd\mu(x)=\pi^{-1}(1+x^2)^{-1}\dd x\)，从而“离零点越近权重越大”被严格化为相位均匀而实轴非均匀。
该构造将非紧对象的压缩实现为一个投影门：将对象提升到相位纤维并在末端以单位圆规范化读出。
与折叠算子 \(\Fold_\infty\) 的横截面投影同型，可得相位投影门的常数预算（推论 \ref{cor:phase-projection-gate-budget}）：任意复合表达式可等价重写为仅在末端调用一次相位投影门的过程层实现。
若要求信息守恒，则必须外置被投影抹去的尺度残差为额外正交轴（命题 \ref{prop:phase-residual-axis}）；该轴的无上界增长在 Hilbert 表示中体现为新增正交自由度，而非“无穷远处含无穷信息”。
在此口径下，单位圆扭曲不再是外参扫描，而是相位坐标化后的内禀位置变量（命题 \ref{prop:phase-intrinsic-twist}）；相位–幅度分离因而成为“紧相位读出/残差外置”的普适输出形式。
压缩发生在相位投影门，守恒发生在残差外置轴，RH 型谱几何发生在紧相位上的酉化判据（命题 \ref{prop:phase-unitary-reduction}）。

\subsubsection{可审计输出：双残差预算、分位证书与端点常数}\label{subsubsec:unit-circle-phase-auditable-outputs}
本段给出三条把相位端点线性化（附录 \ref{app:unit-circle-phase-gate}）、折叠多重度尾部证书（\S\ref{subsec:pom-quantile-budget-law}）与跨分辨率规范差接口（命题 \ref{prop:fold-gauge-anomaly-ldp}）同时闭合的终点断言；其全部输入量均来自既有可审计表格或闭式接口。

\begin{theorem}[结论 A：双残差预算下界与可逆实现的二分相图]\label{thm:unit-circle-phase-double-residual-budget}
在单位圆面积相位坐标门中，伸缩 \(x\mapsto mx\) 由 Möbius 自同构 \(u\mapsto\psi_m(u)\) 诱导，并在 Cayley 半平面坐标 \(\mathfrak{w}(u)=\frac{1+u}{1-u}\) 下满足严格线性化恒等式
\(
\mathfrak{w}(\psi_m(u))=\mathfrak{w}(u)/m
\)
（命题 \ref{prop:scale-linearization-identity-multipliers}，式 \eqref{eq:scale-linearization-identity}），从而尺度链的非紧性被压缩为端点 \(\mathfrak{w}\to 0\) 的精度势（式 \eqref{eq:endpoint-busemann-potential}）。
该紧化门本身可逆；任何不可逆性均来自后续投影/截断并必须按“扩展保存”外置为正交残差轴（公理 \ref{axiom:unit-circle-phase-extension-preserving}；亦见命题 \ref{prop:phase-residual-axis}）。

设同一条流水线同时要求：
\begin{enumerate}
  \item 以 prime-register（CRT）外置纤维标签，使规范投影 \(P_Z\) 在失败概率 \(\le\varepsilon\) 下可逆；
  \item 以局部原子重写回放外置写回支持，使跨分辨率降阶门的可逆化在失败概率 \(\le\varepsilon\) 下成立（其必然触发规范差 \(G_m\)，定义 \ref{def:fold-gauge-anomaly}）。
\end{enumerate}
则任意外置实现都必须在两条互不替代的残差轴上分别满足下界（取 nats 口径）：
\begin{itemize}
  \item \textbf{纤维标签预算（值纤维残差）} 令 \(\gamma_m(\varepsilon)\) 为尾部分位（定义 \ref{def:pom-gamma-quantile}），并记
  \(B_m(\varepsilon):=\exp(m\gamma_m(\varepsilon))\)。
  在区间读出/CRT 外置协议（定义 \ref{def:pom-order-elimination-rewrite}）下，为将“标签超界导致不可唯一重建”的失败概率压到 \(\le\varepsilon\)，prime-register 的模积预算必须满足
  \[
  \boxed{\ \log P_m\ \ge\ \log 2+m\gamma_m(\varepsilon)\ }
  \qquad(\text{等价于 }P_m>2B_m(\varepsilon)),
  \]
  其在该协议意义下为锐阈值（推论 \ref{cor:pom-prime-register-quantile-budget} 及其后“准充要性”段落）。

  \item \textbf{写回支持预算（规范差残差）} 若将可逆化限制为由局部原子重写步骤回放实现，则每一步最多触及 \(3\) 个坐标；因此对任意输入 \(\omega\in\Omega_{m+1}\) 都有纯组合下界
  \[
  \boxed{\ T_{m+1}^{\mathsf S}(\omega)\ \ge\ \left\lceil \frac{G_m(\omega)}{3}\right\rceil\ }
  \]
  （推论 \ref{cor:pom-prime-residual-linear-lb}）。
  在均匀基线 \(\omega\sim\mathrm{Unif}(\Omega_{m+1})\) 下，残差步数的期望速率具有不可约线性下界
  \[
  \liminf_{m\to\infty}\frac{1}{m}\,\mathbb E\!\left[T_{m+1}^{\mathsf S}(\omega)\right]\ \ge\ \frac{4}{27}
  \]
  （同推论 \ref{cor:pom-prime-residual-linear-lb}）。
\end{itemize}
进一步地，在“概率失败 \(\le\varepsilon\)”的指数口径下，第二条可由规范差的速率函数显式定标：
令 \(\beta_m(\varepsilon):=\frac{1}{m}\log(1/\varepsilon)\)，并由命题 \ref{prop:fold-gauge-anomaly-ldp} 定义
\[
I(a)=\sup_{\theta\in\RR}(a\theta-P_G(\theta)),\qquad P_G(\theta)=\log\rho(A_\theta),
\]
再令
\[
\alpha_\varepsilon:=\inf\bigl\{\alpha\in[0,1]: I(\alpha)\ge \beta_m(\varepsilon)\bigr\}.
\]
则任何“每步最多触及 \(3\) 个坐标”的局部回放预算若希望把失败概率压到 \(\le\varepsilon\)，其残差步数必须满足指数尺度下界
\[
\boxed{\ T_{m}\ \gtrsim\ \frac{m}{3}\,\alpha_\varepsilon\ }.
\]
该定理给出一个可审计的二分相图：\(\log P_m\) 由尾部分位 \(\gamma_m(\varepsilon)\) 定标，写回步数由 \(I(\alpha)\) 定标；两者均不被相位门本身消去，因为相位门仅将非紧性挤压到端点精度势 \(\mathfrak{w}\to 0\)，而不改变投影门导致的多对一结构与跨分辨率规范差接口。
\end{theorem}

\begin{theorem}[结论 B：\texorpdfstring{$\gamma_m(\varepsilon)$}{gamma_m(eps)} 的证书化双边界——仅由已核验的 \texorpdfstring{$(r_q,r_{2q})$}{(rq,r2q)} 给出]\label{thm:unit-circle-phase-gamma-two-sided-certificate}
令 \(\beta_m(\varepsilon)=\frac{1}{m}\log(1/\varepsilon)\)，并在均匀微态基线下取输出 \(X\sim w_m\)。
则尾部分位强度 \(\gamma_m(\varepsilon)\)（定义 \ref{def:pom-gamma-quantile}）满足指数尺度双边证书夹逼
\[
\sup_{\substack{q\in\ZZ_{\ge 2}\\ (r_q,r_{2q})\ \mathrm{已审计}\\ \delta(q)<\beta_m(\varepsilon)}} \gamma_q
\ \le\
\gamma_m(\varepsilon)
\ \le\
\inf_{q\in\ZZ_{\ge 2}}
\frac{\log r_q-\log 2+\beta_m(\varepsilon)}{q-1},
\]
其中
\[
\gamma_q:=\frac{\log r_q-\log\varphi}{q},\qquad
\underline f(q):=2\log r_q-\log r_{2q},\qquad
\delta(q):=\log\varphi-\underline f(q)=\log\varphi+\log r_{2q}-2\log r_q.
\]
右侧为 Chernoff--对偶上界（命题 \ref{prop:pom-quantile-chernoff} 及其对偶封装）。
左侧来自 Paley--Zygmund 下包络的分位化推论：当 \(\delta(q)<\beta_m(\varepsilon)\) 时，高多重度事件在指数尺度上的质量严格大于 \(\varepsilon\)，从而 \((1-\varepsilon)\)-尾部分位必不小于对应的 \(\gamma_q\)；
当前审计窗口的具体下包络点 \((\gamma_q,\underline f(q),\delta(q))\) 见表 \ref{tab:fold_collision_multifractal_envelope_pz_q10}（由表格 \(\{r_q\}\) 直接派生，无需任何新增 DP）。
\end{theorem}

\begin{proposition}[结论 C：熵隙常数 \texorpdfstring{$c_\star=\log(2/\varphi)$}{c*=log(2/phi)} 的三重同一性与端点普适指数]\label{prop:unit-circle-phase-entropy-gap-constant}
令
\(
c_\star:=\log(2/\varphi)
\)。
则在本文既定接口下，常数 \(c_\star\) 具有以下三重同一性（均为可审计断言）：
\begin{enumerate}
  \item \textbf{折叠压缩率指数：}平均纤维增长率满足
  \[
  \frac{|\Omega_m|}{|X_m|}=\frac{2^m}{F_{m+2}}=\exp\!\bigl(m c_\star+O(1)\bigr)
  \]
  （见附录 \ref{app:fold-multiplicity} 或命题 \ref{prop:pom-projection-entropy}），从而 \(c_\star\) 是规范投影 \(P_Z\) 的几何压缩率指数。

  \item \textbf{规范差端点指数：}在均匀基线下，规范差比例 \(G_m/m\) 的大偏差速率函数 \(I\) 满足端点闭式
  \[
  I(0)=c_\star
  \]
  （命题 \ref{prop:fold-gauge-anomaly-ldp}），即“完全无需写回”（\(G_m=0\)）是指数稀有事件，其稀有度由同一常数定标。

  \item \textbf{相位端点普适性：}单位圆面积相位门把 \(x\to\infty\) 端点化为 \(\mathfrak{w}\to 0\)，并把伸缩链严格线性化为 \(\mathfrak{w}\mapsto \mathfrak{w}/m\)（命题 \ref{prop:scale-linearization-identity-multipliers}）。
  因此“非紧性”在坐标层面被集中到端点精度势 \(y=-\log|\mathfrak{w}|\) 的加性账本上（式 \eqref{eq:endpoint-busemann-addition-law}），但 \(c_\star\) 作为折叠压缩与规范差端点代价的共同指数仍不可消去。
\end{enumerate}
因此 \(c_\star\) 提供一个跨接口的桥接不变量：它同时刻画“投影导致的平均纤维增长”与“跨分辨率零写回的端点代价”，并为将“相位端点几何”与“投影词成本模型”统一定标提供了共同尺度。
\end{proposition}
